\section{核心功能设计与实现}

\subsection{教务信息管理}
江财教务采用真实校方会话。Android 9 及以上在独立环境中打开官方门户，由用户完成校方登录。应用将校方登录和普通同步拆为两条入口。首次绑定或重连时，账号中心打开交互认证窗口，读取过程先核验本人课表和对应学期，再发布绑定状态。之后的常规查询由独立后台服务复用已有会话，后台读取不会把认证界面推到前台。读取课表时解析校方个人课表网格，转换为包含课程名称、星期、节次范围、周次和地点的统一课程时段。每条记录按上课时段解释，不将同名课程的不同安排合并。周次条件决定条目是否落在选中教学周，起止节次决定其当天的时间范围，校历再把教学周及星期映射为实际日期。这让周视图与日视图可以使用同一批条目。经核验的学期和不支持的结构分别处理，解析失败时保留原缓存。

同步控制依次判断绑定、缓存状态和本次是否需要请求，流程如图\ref{fig:jufe}所示。一般刷新使用六小时新鲜度，重复请求有三十秒限制；手动刷新可跳过新鲜度判断，仍受重复请求限制。页面优先读取当前账号当前学期缓存，新的数据通过校验才替换。本地条目按账号和学期读取。网络错误或解析失败时继续显示已有可用记录，发现账号不匹配则暂停旧缓存的展示。调度区分等待核验与已确认失效，后者持久化为待重连状态，页面恢复也不会重新触发校方认证。

成绩有效记录来自校方相应查询选项。“入学以来”用于整理有记录的学期目录，选择某学期后复核返回报表的学期标识，结果暂存在内存中。成绩、学分和绩点使用校方字段。考试安排沿用校方开放的学期和批次，不通过改写参数扩大范围。当前证据支持空报表的识别，尚缺非空考试报表的实读复核。

\begin{figure}[htbp]
\centering\begin{tikzpicture}
\node[block,text width=37mm] (open) at (-5.2,0) {进入教务或手动刷新};
\node[block,text width=40mm] (cache) at (0,0) {回读同账号、同学期\\最近可用课表};
\node[block,text width=39mm] (decide) at (5.2,0) {检查会话与刷新条件};
\node[block,text width=39mm] (read) at (5.2,-2.0) {后台读取个人课表};
\node[block,text width=40mm] (check) at (0,-2.0) {核验账号、学期与结构};
\node[block,text width=37mm] (save) at (-5.2,-2.0) {保存结果并更新页面};
\node[block,text width=62mm] (keep) at (0,-4.15) {保留可用课程与同步状态\\确认会话失效时提示主动重连};
\draw[flow] (open) -- (cache);
\draw[flow] (cache) -- (decide);
\draw[flow] (decide) -- node[note,right] {允许刷新} (read);
\draw[flow] (read) -- (check);
\draw[flow] (check) -- node[note,above] {通过} (save);
\draw[flow] (check) -- node[note,left] {读取或校验失败} (keep);
\draw[flow] (decide.east) -- ++(3.5mm,0) |- node[note,pos=.84,above] {无需刷新或会话失效} (keep.east);
\end{tikzpicture}

\caption{课表缓存回读与静默同步流程}\label{fig:jufe}
\end{figure}

\subsection{学习通任务聚合}
同名课程可能分属不同教学班，扫描和缓存以课程与教学班标识共同区分。一次扫描先得到教学班列表，再对作业、考试与活动分项查询，记录各自是否完整。聚合时把活动中重复出现且已在作业分项得到的记录去重，再挂接所属教学班及展示所需字段。归并后的任务保留类型、状态、发布或截止信息及可用链接。列表按紧急程度或课程分组，也可按类型及完成状态过滤。例如待批阅作业保留在已完成一侧，缺少可靠绝对截止时间时保留平台提供的可用表达。分页不完整或分项查询失败时，已有有效分项可继续显示。

课程活跃度处理在扫描结果之后，如图\ref{fig:study}所示。设计时选择通知和作业作为观察依据，稳定标识用于区分新记录与旧记录重复出现。确认无更新需要同时具备可用的通知和作业证据。可靠的发布时间用于计算最近活动，稳定的新标识可以补充缺少发布时间的完整记录。首次确认完整但为空的列表从首次可靠观察起算。来源未读全时保留旧状态，避免因查询失败错误隐藏整门教学班。自动隐藏仅影响默认展示。隐藏或活跃度未知的教学班到达十二小时检查间隔后，由同一扫描流程复核。可靠的新作业或通知会撤销自动隐藏，手动隐藏仍由用户解除。手动隐藏和始终显示偏好随账号缓存保存，继续由用户管理。

\begin{figure}[htbp]
\centering\begin{tikzpicture}
\node[block,text width=37mm] (scan) at (-5.2,0) {有效会话\\扫描教学班与分项任务};
\node[block,text width=40mm] (merge) at (0,0) {解析、去重、归并\\失败分项保留旧结果};
\node[block,text width=39mm] (observe) at (5.2,0) {检查作业和通知\\时间、标识及完整性};
\node[block,text width=37mm] (view) at (-5.2,-2.1) {按类型、状态或课程\\筛选并展示可见任务};
\node[block,text width=40mm] (cache) at (0,-2.1) {按账号保存任务及\\活动证据和可见状态};
\node[block,text width=39mm] (state) at (5.2,-2.1) {超期隐藏，有新内容恢复\\证据不足时保留原状态};
\draw[flow] (scan) -- (merge);
\draw[flow] (merge) -- (observe);
\draw[flow] (observe) -- (state);
\draw[flow] (state) -- (cache);
\draw[flow] (cache) -- (view);
\draw[flow] (cache.south) -- (0,-3.45) -- node[note,midway,below] {到期后隐藏项也参与复查} (-7.4,-3.45) |- (scan.west);
\end{tikzpicture}

\caption{学习通任务聚合与课程活跃度的处理顺序}\label{fig:study}
\end{figure}

\subsection{本地待办管理}
待办用本机唯一标识定位记录，保存标题、备注、截止日期、完成标记和创建、更新及完成时间。用户先输入标题建立未完成记录，需要时打开详情补录备注及截止日期。新增和编辑共用字段校验，只有整份记录写入成功才更新界面。勾选完成会写入完成时间，取消勾选则清除这一时间。删除要求二次确认，重启后从应用私有 WebView 存储恢复。列表中的未完成、已完成和全部是同一数据集的筛选结果。待办独立于教务、学习和生活的登录会话，失去外部网络连接时仍可完成新增、编辑和筛选。这为现场演示提供了一条不依赖第三方服务的完整流程。

\subsection{校园生活服务}
酱菜通参照 life-798 的双账号分工和验证码流程，为798设备账号和积分账号分别管理登录通道与会话。设备通道由本人完成验证码登录。账户会话建立后，余额、收藏设备和消费账单分别发起查询并显示平台返回结果。0.8.9 的记录覆盖了这三项数据，并在冷启动后确认会话恢复及再次读取。设备开始与结束需要用户明确触发，连接结果不确定时优先核查已有状态。积分通道具备余额、钱包、任务、签到及兑换流程的代码，实际积分登录、兑换和启停交易仍有单独的验收条件。用户未在阶段验证中发起这些操作。饮水模块使用 jsQR 解析设备二维码，其 Apache 2.0 许可文件随项目保留。

住理生活在生活页单独显示麦庐园项目。官方 App 可按楼栋、楼层和房间选择设备。酱菜通已实现列表与连接状态的界面和逻辑，缺少获准的真实原生数据源时展示当前状态，并提供官方 App 入口。原生洗浴连接、余额、消费记录和设备控制均未在真实环境中验证。

\subsection{统一账号管理}
整合前，用户可能从多个业务页分别寻找登录按钮，一次会话失效也容易在不同页面收到重复提示。统一后的账号中心集中处理绑定、重连与退出状态，继续由各平台自行认证。教务在校方门户完成授权，学习通复用其账号会话，798设备和积分会话可分别管理，住理生活维持独立的服务状态。对学习通和课表，账号校验也参与缓存回读和更新过程，降低切换账号后复用旧数据的可能。业务页保留查询入口，进入重连时转到账号中心的相应服务。这样课表即便保留了旧缓存，也能清楚地提示需要重新认证，用户可以继续读课表并择时处理会话。
